\subsection{ADR-001: estilo de arquitectura}

\textbf{Decisión.} Microservicios por dominio: Usuarios, Cronograma, Solicitudes y documentos, Evaluación y actas, y Notificaciones. Cada servicio tiene su propia base de datos, la aplicación web entra por un API Gateway que valida el JWT y los servicios se avisan los cambios por una cola de mensajes. Los PDF van en un almacén de archivos aparte. Las notas y el acta quedan en el mismo servicio, para que nunca se genere un acta con una nota que no está registrada (riesgo R02).

\textbf{Por qué.} No distribuimos por rendimiento: el sistema tiene pocas decenas de usuarios a la vez, y en una sola máquina un monolito rinde mejor que su versión en microservicios~\cite{Blinowski2022}. Distribuimos porque las partes del proceso no pueden caer juntas: en la semana de predefensas, si fallan la entrega de documentos o los avisos, el tribunal tiene que seguir calificando (RNF-02). Cada dato pide una política distinta: las notas siempre exactas y los avisos con algo de retraso. Y E2 y E3 exigen microservicios y una base distribuida.

\textbf{Opciones descartadas.} El monolito, porque si se cae se cae todo y no sirve para E2 ni E3. El monolito modular, porque sigue siendo un solo proceso con una sola base.

\textbf{Consecuencias.} Hay más piezas que instalar y vigilar (gateway, cola, almacén de archivos y cinco bases), los servicios necesitan contratos claros entre sí y lo que viaja por la cola llega con algo de retraso. Con una sola instancia, el API Gateway es un punto único de falla (riesgo R03).

El registro completo está en \nolinkurl{docs/adr/ADR-001-estilo-arquitectura.md}.
